\documentclass[11pt]{article}
\usepackage{amsmath}
\usepackage[hidelinks]{hyperref}
\title{Earning a formation notice by moving through a sensor window}
\date{27 November 2026}
\begin{document}
\maketitle
\paragraph{Audit.}
Use two adjacent ordered blocks of identical
acceleration- and speed-bounded robots, all initially and finally
at rest. The center starts at coordinate zero. The demanded block
is stored as a current disposition on a fixed beacon, not known
to the moving agents. Only the center needs the result here:
every noncentral robot in each block follows a preinstalled
rigid translation outwards. The hard gap slack per gap is
at least \(2B(T)\), where \(B\) is the rest-to-rest displacement
budget of PR~709. This makes both wide directions lawful; no
claim is made for a tighter link budget. Receiving a typed result,
being geometrically within the read strip, and spending radio
allowance are distinct events. The motion and beacon are authored
one-dimensional plants; the selected SubT broker and its
application callbacks are executed.

\paragraph{Exact instantaneous acquisition boundary.}
A functioning instantaneous read is possible only upon reaching
coordinate \(p>0\) (a further upper bound at least \(p\) is
allowed). Absent a read the center cannot distinguish the
two demands. Write \(P=B(T)\) and
\(\tau_* = B^{-1}(p)\). If a gain greater than \(P\) is possible,
the adverse world must visit \(p\) and later finish at or left
of \(-z\), where \(z=\delta-P>0\). Such a continuous rest-ended
trajectory has a maximum position \(q\ge p\) with zero velocity
at an interior time \(t_*\). Rest-to-rest bounds give
\(q\le B(t_*)\) and \(q+z\le B(T-t_*)\). Therefore \(t_*\ge\tau_*\)
and \(z\le B(T-\tau_*)-p\). Conversely, go to \(p\) at rest at
\(\tau_*\), obtain the current sign, then move to \(+z\) or
\(-z\). Both require at most \(p+z\) of remaining travel.
The exact maximum common target is consequently
\[
             P+\max(0,B(T-\tau_*)-p),
\]
using the blind motion if \(\tau_*>T\) or the additional
term is nonpositive. Even allowing the probe to be crossed in
motion cannot improve this worst-case bound: the adverse
continuation must turn. This is not true for a cue that is
delivered without the position requirement.

For the construction let all outer robots translate outwards
under the full envelope \(F_T(t)\). Every center trajectory from
and to rest has instantaneous speed bounded by that same full
envelope, so \(|x(t)|\le F_T(t)\). Each gap adjacent to the
center gains \(F_T(t)\pm x(t)\), in \([0,2P]\);
the other gaps are unchanged. This gives a continuous,
collision-clear, geometric-link-preserving witness. There is
assumed to be adequate free corridor extent.

\paragraph{Dwell can make motion during acquisition essential.}
Now require a continuous unit of exposure in the strip
\([3/4,1]\); after that the beacon may reply. Fix acceleration
2, speed cap at least 3, and deadline \(T=3\). The common probe
goes from 0 to 1 at rest by \(3/2\), with acceleration
magnitude \(16/9\). At time 1 it is at \(7/9\).
It then moves left with acceleration 2, reaching \(3/4\)
at time 2 with velocity \(-1\).
The past path during \([1,2]\) is entirely in the strip,
as shown by its exact extrema and monotonic segments.
Only at the achieved read cut is the stored demand used.

At that common moving state, the negative branch finishes at
\(-1/8\) at time 3 on the rest-to-rest envelope launched
from 1 at time \(3/2\). The positive branch brakes immediately,
stops at \(1/2\), and waits. The exact demonstrated target is
\(37/8\), above the blind guarantee \(P=9/2\).
If the result is unavailable, the center uses a preinstalled
\(4{:}1\) convex combination of the two suffixes. It ends at
zero, at rest, restoring the blind \(9/2\) floor even after the
probe has moved. This costs motion and is not an information
gain under loss.

If acquisition instead demands that the center be stationary
at a single point during the full exposure, the adverse branch
needs at least \(\sqrt{3/2}+1+\sqrt{7/4}>3\) time: first
move at rest to a valid point beyond \(3/4\), dwell, and then
end at or below \(-1/8\). So stopping for the read fails in
this family. The positive moving example is not the optimum
for general dwell geometry.

\paragraph{Executed evidence path.}
At tick 2 the authored beacon receives the qualified past
geometric certificate. Only then can it emit the current
typed two-byte sign. The center branches only on a current
bound ROS callback from that beacon. If the strip is narrowed
so that the same motion leaves it, there is no certificate.
A stale response can physically arrive without authorizing
the branch. An erased response never does. A prior four-byte
report from the same five-byte sender at time 1 still occupies
the native inclusive window at time 2 and blocks the response.
The same report at time \(7/8\) no longer does. The earlier
old report is actually delivered in those controls; we do not
cancel that work to gain capacity. The radio path quality and
the instant application timing at this cut are declared
test seams, not measured optics or packet time.

The selected source replay is matched to the authored geometric
and strict budget predictions in 48 histories, with 4 qualified
actuations and 44 neutral fallback outcomes. The continuous
proof, independent dyadic path check and negative controls
do not supply a physical sensor deployment. An equally
informed direct controller with the same evidence and resources
ties. There is no independent mission or R6 acceptance claim.
\end{document}
